\section{Рабочий проект}

В процессе разработки программно-информационной поисковой системы были реализованы два самостоятельных программных компонента: библиотека регулярных выражений на языке Си и поисковое приложение на C++ с использованием фреймворка Qt.

\subsection{Библиотека регулярных выражений}

\subsubsection{Состав библиотеки и условия применения}

Библиотека регулярных выражений поставляется в виде двух файлов:
\verb`regex-lib.c` -- реализация всех внутренних и экспортируемых
функций, и \verb`regex-lib.h` -- заголовочный файл с объявлениями
экспортируемых функций, непрозрачного типа \verb`struct re_chain`
и перечисления кодов ошибок \verb`enum re_error_type`. Для подключения
библиотеки к проекту достаточно добавить оба файла в систему сборки
и включить заголовочный файл. В C++-проектах подключение выполняется
без дополнительных настроек благодаря блоку \verb`extern "C"`
в заголовочном файле.

Библиотека протестирована на следующих платформах и компиляторах:
\begin{itemize}
	\item Windows~10/11, Clang~17;
	\item Linux (Ubuntu~22.04), GCC~12;
	\item онлайн-компиляторы на базе GCC и Clang (OnlineGDB и аналоги).
\end{itemize}

При компиляции не возникает предупреждений. Проверка санитайзерами
адресов и неопределённого поведения (\verb`-fsanitize=address,undefined`)
не выявила ошибок.

\subsubsection{Сборка библиотеки}

Для сборки библиотеки не требуется система автоматической сборки
и сторонние зависимости -- достаточно компилятора, поддерживающего
стандарт C99 или новее. Команды сборки представлены
в таблице~\ref{tab:build_commands}.

\begin{xltabular}{\textwidth}{|C{3.5cm}|T|}
	\caption{Команды сборки библиотеки}
	\label{tab:build_commands} \\
	\hline
	\thead{Режим} & \thead{Команда} \\
	\hline
	\finishhead
	Отладочная сборка тестов & \verb`gcc -g -Wall -Wextra -fsanitize=address,undefined`
	\verb`regex-tests.c -o regex-tests-dbg` \\
	\hline
	Оптимизированная сборка тестов & \verb`gcc -O2 regex-tests.c -o regex-tests` \\
	\hline
	Релизная сборка библиотеки & \verb`gcc -O2 -DNDEBUG -c regex-lib.c -o regex-lib.o` \verb`gcc program.c regex-lib.o -o program` \\
\end{xltabular}

Команды идентичны при использовании Clang: достаточно заменить
\verb`gcc` на \verb`clang`. Тестовый файл \verb`regex-tests.c` включает
реализацию библиотеки напрямую директивой \verb`#include "regex-lib.c"`,
что позволяет собирать тесты единственной командой без раздельной
компиляции. При интеграции библиотеки в сторонний проект рекомендуется
раздельная компиляция.



\subsubsection{Экспортируемые функции библиотеки}

Библиотека предоставляет C API, объявленное в заголовочном файле \verb`regex-lib.h`. Все экспортируемые функции разделены на семейства по выполняемой операции. Для каждого семейства предоставляются две версии: принимающая строку шаблона и выполняющая его компиляцию при каждом вызове, и версия с префиксом \verb`obj_`, принимающая заранее скомпилированный объект типа \verb`struct re_chain*`. Использование \verb`obj_`-версий предпочтительно при многократном применении одного шаблона. 
Помимо экспортируемых функций, заголовочный файл объявляет непрозрачный тип \verb`struct re_chain` и перечисление кодов ошибок \verb`enum re_error_type`, что позволяет потребителю API обрабатывать конкретные коды ошибок без доступа к внутренней реализации.

Дополнительно в файле исходного кода библиотеки доступны две вспомогательные функции, не входящие в публичный API и предназначенные для отладки:
\verb`print_re_blocks` -- форматированный вывод внутреннего представления скомпилированного регулярного выражения в консоль и
\verb`get_error_msg` -- получение текстового описания кода ошибки на английском языке.
\\

\subsubsubsection{Компиляция шаблона регулярного выражения}

Функции этого семейства выполняют разбор строки шаблона и построение внутреннего представления регулярного выражения, а также освобождение занимаемой им памяти. Аргументы функций приведены в таблицах~\ref{tab:api_compile} и~\ref{tab:api_free_re}.

Возвращаемое значение функций компиляции~-- указатель на непрозрачный тип \verb`struct re_chain*`:
\begin{enumerate}
	\item Не-\verb`NULL` -- компиляция успешна, возвращен объект регулярного выражения.
	\item \verb`NULL` -- ошибка в строке шаблона.
\end{enumerate}

Прототипы функций:
\begin{enumerate}
	\item \verb`struct re_chain* compile_re(char* pattern)`
	\item \verb`struct re_chain* compile_re2(char* pattern, int* errcode, int* errpos)`
\end{enumerate}

\begin{xltabular}{\textwidth}{|C{3cm}|T|C{3cm}|}
	\caption{Аргументы функций компиляции}
	\label{tab:api_compile} \\
	\hline
	\thead{Аргумент} & \thead{Описание} & \thead{Допускается\\ NULL?} \\
	\hline
	\finishhead
	\verb`char* pattern` & Строка шаблона регулярного выражения &  \\
	\hline
	\verb`int* errcode` (для \verb`compile_re2`) & Указатель для записи кода ошибки компиляции & Да \\
	\hline
	\verb`int* errpos` (для \verb`compile_re2`) & Указатель для записи позиции символа в шаблоне, при разборе которого обнаружена ошибка & Да \\
	\hline
\end{xltabular}

Функция \verb`free_re` освобождает всю память, выделенную под объект регулярного выражения и все его внутренние структуры. Должна быть вызвана для каждого объекта, созданного функциями \verb`compile_re` или \verb`compile_re2`, после окончания его использования.
Функция ничего не возвращает.

Прототип функции:
\begin{enumerate}
	\item \verb`void free_re(struct re_chain* re_obj)`
\end{enumerate}

\begin{xltabular}{\textwidth}{|C{3cm}|T|C{3cm}|}
	\caption{Аргументы функции освобождения памяти}
	\label{tab:api_free_re} \\
	\hline
	\thead{Аргумент} & \thead{Описание} & \thead{Допускается\\ NULL?} \\
	\hline
	\finishhead
	\verb`struct re_chain*` \verb`re_obj` & Указатель на скомпилированный объект регулярного выражения, память которого должна быть освобождена &  \\
\end{xltabular}

\subsubsubsection{Поиск позиции совпадения}

Функции этого семейства находят первое совпадение шаблона в строке и возвращают его начальную позицию и длину. Аргументы функций приведены в таблице~\ref{tab:api_matchpos}.

Возвращаемое значение -- целое число: 
\begin{enumerate}
	\item >= 0 -- захват успешен, его начальная позиция во входной строке.
	\item -1 -- совпадений во входной строке не найдено.
	\item < -1 -- ошибка в строке шаблона с возвращенным кодом.
\end{enumerate}

Прототипы функций:
\begin{enumerate}
	\item \verb`int matchpos(char* pattern, char* text, int* capture_data)`
	\item \verb`int obj_matchpos(struct re_chain* re_obj, char* text, int* capture_data)`
\end{enumerate}

\begin{xltabular}{\textwidth}{|C{3cm}|T|C{3cm}|}
	\caption{Аргументы функций поиска позиции совпадения}
	\label{tab:api_matchpos} \\
	\hline
	\thead{Аргумент} & \thead{Описание} & \thead{Допускается\\ NULL?} \\
	\hline
	\finishhead
	\verb`char* pattern` / \verb`struct re_chain*` \verb`re_obj` & Строка шаблона регулярного выражения / Откомпилированный объект регулярного выражения &  \\
	\hline
	\verb`char* text` & Входная строка для поиска & \\
	\hline
	\verb`int* capture_data` & Указатель на целое: при вызове -- начальная позиция поиска (если значение > 0); после вызова -- длина найденного совпадения или позиция ошибки & Да \\
\end{xltabular}


\subsubsubsection{Валидация строки}

Функции этого семейства проверяют, содержит ли входная строка хотя бы одно совпадение с шаблоном, и возвращают булево значение. Аргументы функций приведены в таблице~\ref{tab:api_ismatch}.

Возвращаемое значение -- целое число: 
\begin{enumerate}
	\item 1 -- входная строка удовлетворяет регулярному выражению.
	\item 0 -- входная строка не подходит под регулярное выражение.
	\item < -1 -- ошибка в строке шаблона с возвращенным кодом.
\end{enumerate}

Прототипы функций:
\begin{enumerate}
	\item \verb`int is_match(char* pattern, char* text)`
	\item \verb`int obj_is_match(struct re_chain* re_obj, char* text)`
\end{enumerate}

\begin{xltabular}{\textwidth}{|C{3cm}|T|C{3cm}|}
	\caption{Аргументы функций валидации строки}
	\label{tab:api_ismatch} \\
	\hline
	\thead{Аргумент} & \thead{Описание} & \thead{Допускается\\ NULL?} \\
	\hline
	\finishhead
	\verb`char* pattern` / \verb`struct re_chain*` \verb`re_obj` & Строка шаблона регулярного выражения / Скомпилированный объект регулярного выражения &  \\
	\hline
	\verb`char* text` & Входная строка для валидации &  \\
\end{xltabular}

\subsubsubsection{Извлечение одного совпадения с группами захвата}

Функции этого семейства возвращают NULL-терминированный массив строк, содержащий полное совпадение и все захваченные группы, аналогично функции \verb`match()` в языке JavaScript. Аргументы функций приведены в таблице~\ref{tab:api_match}.

Возвращаемое значение -- NULL-терминированный массив malloc-строк: 
\begin{enumerate}
	\item Указатель на указатель на строку -- захват успешен: первый элемент массива -- общий захват, последующие (пока не встретится указатель на NULL) -- группы захвата.
	\item Указатель на указатель на NULL -- входная строка не подходит под регулярное выражение.
	\item NULL -- ошибка в строке шаблона с возвращенным кодом.
\end{enumerate}

Прототипы функций:
\begin{enumerate}
	\item \verb`char** match(char* pattern, char* text)`
	\item \verb`char** obj_match(struct re_chain* re_obj, char* text)`
\end{enumerate}

\begin{xltabular}{\textwidth}{|C{3cm}|T|C{3cm}|}
	\caption{Аргументы функций извлечения одного совпадения}
	\label{tab:api_match} \\
	\hline
	\thead{Аргумент} & \thead{Описание} & \thead{Допускается\\ NULL?} \\
	\hline
	\finishhead
	\verb`char* pattern` / \verb`struct re_chain*` \verb`re_obj` & Строка шаблона регулярного выражения / Скомпилированный объект регулярного выражения &  \\
	\hline
	\verb`char* text` & Входная строка для поиска &  \\
\end{xltabular}

\subsubsubsection{Извлечение всех совпадений}

Функции этого семейства возвращают NULL-терминированный массив строк, содержащий все непересекающиеся совпадения регулярного выражения во входной строке, аналогично функции \verb`match()` с флагом \verb`g` в языке JavaScript. Аргументы функций приведены в таблице~\ref{tab:api_matches}.

Возвращаемое значение -- NULL-терминированный массив malloc-строк:
\begin{enumerate}
	\item Указатель на массив строк -- все найденные совпадения; если совпадений нет, первым элементом массива будет \verb`NULL`.
	\item \verb`NULL` -- ошибка в строке шаблона.
\end{enumerate}

Прототипы функций:
\begin{enumerate}
	\item \verb`char** matches(char* pattern, char* text)`
	\item \verb`char** obj_matches(struct re_chain* re_obj, char* text)`
\end{enumerate}

\begin{xltabular}{\textwidth}{|C{3cm}|T|C{3cm}|}
	\caption{Аргументы функций извлечения всех совпадений}
	\label{tab:api_matches} \\
	\hline
	\thead{Аргумент} & \thead{Описание} & \thead{Допускается\\ NULL?} \\
	\hline
	\finishhead
	\verb`char* pattern` / \verb`struct re_chain*` \verb`re_obj` & Строка шаблона регулярного выражения / Скомпилированный объект регулярного выражения &  \\
	\hline
	\verb`char* text` & Входная строка для поиска &  \\
\end{xltabular}

\subsubsubsection{Замена совпадений в буфере фиксированного размера}

Функции этого семейства выполняют замену совпадений шаблона на строку замены в предоставленном пользователем буфере фиксированного размера. Аргументы функций приведены в таблице~\ref{tab:api_nreplace}.

Возвращаемое значение~-- целое число:
\begin{enumerate}
	\item $\ge 0$ -- число успешно выполненных замен.
	\item $-1$ -- в буфере \verb`dest` закончилось место -- результат неполный.
	\item $< -1$ -- ошибка в строке шаблона с возвращенным кодом.
\end{enumerate}

Прототипы функций:
\begin{enumerate}
	\item \verb`int re_nreplace(char* pattern, char* repl, char* text, char* dest, int n, int limit)`
	\item \verb`int obj_re_nreplace(struct re_chain* re_obj, char* repl, char* text, char* dest,` \verb`int n, int limit)`
\end{enumerate}

\begin{xltabular}{\textwidth}{|C{3cm}|T|C{3cm}|}
	\caption{Аргументы функций замены с фиксированным буфером}
	\label{tab:api_nreplace} \\
	\hline
	\thead{Аргумент} & \thead{Описание} & \thead{Допускается\\ NULL?} \\
	\hline
	\finishhead
	\verb`char* pattern` / \verb`struct re_chain*` \verb`re_obj` & Строка шаблона регулярного выражения / Скомпилированный объект регулярного выражения &  \\
	\hline
	\verb`char* repl` & Строка замены; может содержать ссылки на группы захвата вида \verb`$n` или \verb`\n` (где \textit{n} от 0 до 9) &  \\
	\hline
	\verb`char* text` & Исходная строка, в которой выполняется замена &  \\
	\hline
	\verb`char* dest` & Указатель на буфер для записи результата замены &  \\
	\hline
	\verb`int n` & Размер буфера \verb`dest` в байтах &  \\
	\hline
	\verb`int limit` & Максимальное число замен; 0~-- без ограничений &  \\
\end{xltabular}

\subsubsubsection{Замена совпадений с динамическим выделением памяти}

Функции этого семейства выполняют замену совпадений шаблона на строку замены, автоматически выделяя результирующий буфер необходимого размера в динамической памяти. Аргументы функций приведены в таблице~\ref{tab:api_dupreplace}.

Возвращаемое значение~-- указатель на \verb`malloc`-строку или \verb`NULL`:
\begin{enumerate}
	\item Не-\verb`NULL` -- указатель на строку с результатом замены (должна быть освобождена вызывающей стороной).
	\item \verb`NULL` -- ошибка в строке шаблона (код ошибки при этом записывается по указателю \verb`repl_cnt`).
\end{enumerate}

Прототипы функций:
\begin{enumerate}
	\item \verb`char* re_dupreplace(char* pattern, char* repl, char* text, int limit, int* repl_cnt)`
	\item \verb`char* obj_re_dupreplace(struct re_chain* re_obj, char* repl, char* text, int limit,` \verb`int* repl_cnt)`
\end{enumerate}

\begin{xltabular}{\textwidth}{|C{3cm}|T|C{3cm}|}
	\caption{Аргументы функций замены с динамическим выделением памяти}
	\label{tab:api_dupreplace} \\
	\hline
	\thead{Аргумент} & \thead{Описание} & \thead{Допускается\\ NULL?} \\
	\hline
	\finishhead
	\verb`char* pattern` / \verb`struct re_chain*` \verb`re_obj` & Строка шаблона регулярного выражения / Скомпилированный объект регулярного выражения &  \\
	\hline
	\verb`char* repl` & Строка замены; может содержать ссылки на группы захвата вида \verb`$n` или \verb`\n` (где \textit{n} от 0 до 9) &  \\
	\hline
	\verb`char* text` & Исходная строка, в которой выполняется замена &  \\
	\hline
	\verb`int limit` & Максимальное число замен; 0 -- без ограничений &  \\
	\hline
	\verb`int* repl_cnt` & Указатель для записи числа выполненных замен или отрицательного кода ошибки & Да \\
\end{xltabular}


\subsubsection{Тестирование библиотеки регулярных выражений}

Для верификации корректности библиотеки проведено модульное и интеграционное тестирование. Модульное тестирование охватывает каждое семейство экспортируемых функций в отдельности: тесты модулей \verb`check_correctness`, \verb`matchpos`, \verb`matchstart` проверяют функции компиляции и сопоставления изолированно. Интеграционное тестирование охватывает совместную работу нескольких компонентов: модуль \verb`is_match` проверяет цепочку «компиляция через \verb`compile_re2` -- сопоставление через \verb`obj_is_match`»; модули \verb`matcharr`, \verb`matches` и \verb`replace` проверяют полный конвейер «компиляция -- сопоставление -- извлечение подстрок / замена».

В файле \verb`regex-tests.c` содержится встроенный набор автоматизированных тестов. Тестирование библиотеки выполняется в первых двух профилях сборки (отладочной и оптимизированной), описанных в таблице~\ref{tab:build_commands}.

Профиль отладочной сборки используется для выявления ошибок работы с памятью и неопределённого поведения средствами санитайзеров AddressSanitizer и UndefinedBehaviorSanitizer. Время выполнения тестов в этом профиле не показательно из-за накладных расходов инструментации.

Профиль оптимизированной сборки выполняет тестирование в условиях, приближённых к промышленной эксплуатации и включает в себя агрессивные оптимизации компилятора, которые могут проявить скрытое неопределённое поведение (UB), не видимое при отладочной сборке.

В отладочном профиле вывод санитайзеров должен быть чистым -- любое
сообщение от AddressSanitizer или UBSanitizer означает ошибку в реализации.
В оптимизированном профиле дополнительно фиксируется время выполнения
тест-сьюта, которое при сборке с \verb`-O2` составляет 0,022~с.

Файл \verb`regex-tests.c` включает реализацию библиотеки директивой \verb`#include "regex-lib.c"`, что позволяет обойтись без раздельной компиляции при запуске тестов.

Тесты организованы в модули с помощью макросов \verb`NEW_MODULE(name)` и \verb`FINISH_MODULE()`. Результат отдельного теста учитывается макросами \verb`TEST_OK(...)` и \verb`TEST_FAIL(...)`, фиксирующими успешное и неудачное прохождение теста соответственно. По завершении каждого модуля выводится число пройденных тестов. Тесты, помеченные макросом \verb`SPECIAL_TEST`, не учитываются в статистике -- они проверяют граничные случаи, поведение в которых не стандартизировано и различается среди промышленных библиотек регулярных выражений (PCRE, Go \verb`regexp`, Notepad++).

Тест-сьют разделён на шесть модулей, описанных ниже.

\paragraph{Модуль \texttt{check\_correctness}}

Модуль проверяет корректность диагностики синтаксических ошибок в шаблонах регулярных выражений. Для каждого тестового шаблона функция \verb`assert_check_correctness` вызывает \verb`compile_re2` и сравнивает возвращённый код ошибки с эталонным значением из перечисления \verb`re_error_type`. В выводе для каждого теста указывается код ошибки, её текстовое описание и позиция ошибочного символа в шаблоне.

Тесты охватывают все 16 кодов ошибок, определённых в библиотеке, включая корректные шаблоны (ожидаемый результат \verb`RE_OK`), а также условную компиляцию для режима \verb`STRICT_CHARCLASSES`, запрещающего экранирование букв, не используемых в символьных классах (например \verb`j`, \verb`q`). Всего в модуле 38 тестов.

% Настройка листингов как рисунков (подпись снизу, сквозная нумерация с рисунками)
\renewcommand{\lstlistingname}{Рисунок}
\lstset{
	captionpos=b,                    % подпись снизу
	aboveskip=10pt,                  % отступ сверху (как у рисунков)
	belowskip=0pt,
%	belowcaptionskip=10pt  не работает
}

% Привязка счётчика листингов к разделу (как у рисунков в vkr)
\makeatletter
\@addtoreset{lstlisting}{section}
\makeatother
\renewcommand{\thelstlisting}{\thesection.\arabic{lstlisting}}

\setcounter{lstlisting}{0}

Результат выполнения модуля \verb`check_correctness` представлен на рисунке~\ref{fig:test_correctness}.

\begin{lstlisting}[caption={Результат выполнения модуля check\_correctness}, label={fig:test_correctness}]
	--- Module "check_correctness" ---
	OK 0 "OK"
	OK -100 "Unescaped ^ not in the beginning of regex or group" ('^' at 1)
	OK -99 "Unescaped $ not in the end of regex" ('$' at 0)
	OK -98 "Preceding token is not quantifiable" ('+' at 14)
	OK -98 "Preceding token is not quantifiable" ('*' at 12)
	OK -98 "Preceding token is not quantifiable" ('?' at 3)
	OK -97 "Empty square brackets" (']' at 3)
	OK -97 "Empty square brackets" (']' at 4)
	OK -96 "Missing opening square bracket" (']' at 5)
	OK -95 "Missing closing square bracket" ('' at 8)
	OK -94 "Range with character class is inacceptable" ('-' at 3)
	OK -94 "Range with character class is inacceptable" ('-' at 2)
	OK 0 "OK"
	OK -93 "Missing character after escape character" ('' at 4)
	OK 0 "OK"
	OK -100 "Unescaped ^ not in the beginning of regex or group" ('^' at 5)
	OK 0 "OK"
	OK 0 "OK"
	OK 0 "OK"
	OK 0 "OK"
	OK -92 "Invalid characters inside curly brackets" ('a' at 8)
	OK -92 "Invalid characters inside curly brackets" (',' at 7)
	OK -91 "Quantifier range is out of order" ('}' at 7)
	OK -90 "Quantifier range is too large" ('9' at 14)
	OK 0 "OK"
	OK 0 "OK"
	OK -88 "Missing closing parenthesis" ('' at 3)
	OK -89 "Missing opening parenthesis" (')' at 2)
	OK 0 "OK"
	OK -88 "Missing closing parenthesis" ('' at 15)
	OK -89 "Missing opening parenthesis" (')' at 14)
	OK 0 "OK"
	OK -87 "Character range is out of order" ('y' at 8)
	OK -85 "Invalid regex modifier" ('j' at 3)
	OK -88 "Missing closing parenthesis" ('' at 4)
	OK 0 "OK"
	OK 0 "OK"
	OK -86 "Invalid character class" ('j' at 2)
	All module tests (38) passed
\end{lstlisting}

\paragraph{Модуль \texttt{matchpos}}

Модуль проверяет весь основной синтаксис регулярных выражений на основе функции \verb`matchpos` -- поиска первого совпадения шаблона в строке. Функция \verb`assert_match` вызывает \verb`matchpos` с нулевой начальной позицией и сравнивает найденную подстроку с эталонной строкой. Для тестов, ожидающих отсутствие совпадений, эталонным значением является \verb`NULL`.

Тесты охватывают следующие возможности библиотеки:
\begin{itemize}
	\item базовые метасимволы и квантификаторы (\verb`.`, \verb`*`, \verb`+`, \verb`?`, \verb`{n,m}`);
	\item наборы символов (\verb`[...]`, \verb`[^...]`) и диапазоны;
	\item якоря \verb`^` и \verb`$` в том числе в составе альтераций;
	\item альтерацию с выбором наиболее раннего и наиболее длинного захвата;
	\item группы захвата, незахватываемые группы \verb`(?:...)` и обратные ссылки \verb`\1`;
	\item символьные классы \verb`\d`, \verb`\w`, \verb`\s`, \verb`\a`, \verb`\p` и их инверсии;
	\item ленивые квантификаторы \verb`*?`, \verb`+?`, \verb`??`;
	\item модификаторы \verb`(?s)`, \verb`(?m)`, \verb`(?i)`, \verb`(?I)` (в том числе регистронезависимый поиск с кириллицей в UTF-8);
	\item работу с пустыми шаблонами и пустыми строками. 
\end{itemize}

Три теста помечены как особые (\verb`SPECIAL_TEST`): выбор между альтерациями с одинаковой начальной позицией и ленивый захват внутри групп. Всего в модуле 136 тестов, из которых 3 особых.

Результат выполнения модуля \verb`matchpos` представлен на рисунке~\ref{fig:test_matchpos}.

\begin{lstlisting}[caption={Результат выполнения модуля matchpos}, label={fig:test_matchpos}]
	--- Module "matchpos" ---
	OK "abc"
	OK "dksst"
	OK "bfd"
	OK "a123bc9"
	OK "heo"
	OK NULL
	OK "abcktrrr"
	OK "b 4d_"
	OK NULL
	OK ""
	OK ""
	OK ""
	OK ""
	OK ""
	OK NULL
	OK "aaaabw"
	OK "test   string 12"
	OK "Helo"
	OK "привет"
	OK "at"
	OK "[][q[w"
	OK "aaabab"
	OK "c"
	OK "1111"
	OK "aa"
	OK ""
	OK "aaaaa"
	OK ""
	OK "aaaa"
	OK "aba"
	OK NULL
	OK "1234"
	OK "3456"
	OK NULL
	OK ""
	OK ""
	OK "27:15"
	OK "2:15"
	OK "1:27"
	OK NULL
	OK "https://swsu.ru"
	OK NULL
	OK NULL
	OK "http://www.swsu.ru"
	OK "abc@af-af.de"
	OK NULL
	OK NULL
	OK NULL
	OK "f"
	OK "wq"
	OK NULL
	OK NULL
	OK "def"
	OK "abc"
	OK ""
	OK "bddd"
	OK "1st"
	OK "stecond"
	OK NULL
	OK "ddd"
	OK NULL
	* OK "aabcc"
	OK "abcd a1 asdd asd asad da"
	OK "abcd a1 asdd asd asa"
	OK "abcd a1 asdd asd asad das"
	OK NULL
	OK "abcdef"
	OK "abcabc"
	OK ""
	OK "1abcabc"
	OK "qabqwcdab"
	OK "aaaaaaaa"
	OK "aaaaaaaaaaaa"
	OK "abcdet2b!end"
	OK NULL
	OK "abcdet2end"
	OK "cd"
	OK ""
	OK ""
	OK "https://www.youtube.com/"
	OK "https://youtube-nocookie.com/"
	OK "http://youtu.be/"
	OK NULL
	OK NULL
	OK "https://i.ytimg.com/"
	OK "abed"
	OK NULL
	OK "abc"
	OK NULL
	OK "134"
	OK "1345"
	OK "13445"
	OK "124"
	OK "abABC!"
	OK "abAB\|_"
	OK "abc45"
	OK "abc_d1"
	OK "ab#12$YZA"
	OK "abc,        d!ef"
	OK "a   b
	"
	OK "ba"
	OK NULL
	OK "ab"
	OK "a"
	OK "abcc"
	OK "a"
	OK ""witch" and her "broom""
	OK ""witch""
	OK "<p> <b> Wikipedia </b> encyclopedia. </p>"
	OK "<p>"
	OK "aa"
	OK "aaa b"
	OK "{FW}"
	OK "869759002514931"
	OK "3.1415926"
	OK "y!y!"
	OK NULL
	OK "#y!y!"
	OK NULL
	OK "abab"
	OK "abab"
	OK "aa"
	OK NULL
	OK ""She's the one!""
	OK "_'a b c' 'd e f' !"
	* FAIL NULL != ""She's the one!""
	* FAIL [] != ["  test s   ", "test s"]
	OK "line1"
	OK ".*?"
	OK "line1
	line2"
	OK "b"
	OK "b"
	OK "bc"
	OK "TEst"
	OK "TEst"
	OK "ATest"
	OK "поддержка"
	OK "поддержка"
	OK "Тестовый поиск"
	All module tests (136) passed
\end{lstlisting}

\paragraph{Модуль \texttt{matchstart}}

Модуль проверяет корректность поиска с заданной начальной позицией -- передачи положительного значения по указателю \verb`capture_data` при вызове \verb`matchpos`. Функция \verb`assert_matchstart` принимает дополнительный параметр \verb`start_pos` и вызывает \verb`matchpos` с этим значением.

Тесты проверяют: поиск с ненулевой позиции, взаимодействие начальной позиции с якорем \verb`^` (поиск не должен находить совпадение, если якорь не применим с заданной позиции), а также поиск с позицией, выходящей за пределы строки. Всего 4 теста.

Результат выполнения модуля \verb`matchstart` представлен на рисунке~\ref{fig:test_matchstart}.
\\

\begin{lstlisting}[caption={Результат выполнения модуля matchstart}, label={fig:test_matchstart}]
	--- Module "matchstart" ---
	OK " b"
	OK "</b>"
	OK NULL
	OK NULL
	All module tests (4) passed
\end{lstlisting}

\paragraph{Модуль \texttt{is\_match}}

Модуль проверяет корректность функции \verb`obj_is_match` -- валидации строки на соответствие шаблону. Функция \verb`assert_is_match` для каждого теста компилирует шаблон через \verb`compile_re2`, вызывает \verb`obj_is_match` и сравнивает булев результат с эталонным. Использование \verb`obj_`-версии дополнительно проверяет корректность работы с предварительно скомпилированным объектом.

Тесты охватывают: базовое совпадение и его отсутствие, якоря в составе групп, а также практический шаблон проверки на простоту чисел в унарной системе счисления \verb`^1?$|^(11+)\1+$` от 1 до 17, и шаблон валидации адреса электронной почты, используемый в поле HTML \verb`<input type="email">`. Всего 20 тестов.

Результат выполнения модуля \verb`is_match` представлен на рисунке~\ref{fig:test_ismatch}.

\begin{lstlisting}[caption={Результат выполнения модуля ismatch}, label={fig:test_ismatch}]
	--- Module "is_match" ---
	OK true
	OK false
	OK true
	OK false
	OK true
	OK false
	OK true
	OK true
	OK true
	OK true
	OK true
	OK false
	OK true
	OK false
	OK true
	OK false
	OK true
	OK false
	OK true
	OK false
	All module tests (20) passed
\end{lstlisting}

\paragraph{Модуль \texttt{matcharr}}

Модуль проверяет корректность функций \verb`match` и \verb`matches` через функции-обёртки \verb`assert_matcharr` и \verb`assert_matches` соответственно. Обе функции сравнивают возвращённый NULL-терминированный массив строк с эталонным массивом, переданным через вариадический макрос \verb`ARR(...)`.

Тесты функции \verb`match` проверяют: возврат полного совпадения без групп, возврат групп захвата при простых и вложенных группах, незахватываемые группы \verb`(?:...)`, взаимодействие групп с квантификаторами, а также сложный вложенный шаблон с несколькими уровнями альтераций. Всего 14 тестов.

Результат выполнения модуля \verb`matcharr` представлен на рисунке~\ref{fig:test_matcharr}.

\begin{lstlisting}[caption={Результат выполнения модуля matcharr}, label={fig:test_matcharr}]
	--- Module "matcharr" ---
	OK ["x"]
	OK ["x", "x"]
	OK ["abcdef", "abc", "def"]
	OK ["abcdef", "def"]
	OK ["abcdef", "abc"]
	OK ["aaaaaaaa", "aaaa", "aa"]
	OK ["aaaaaaaa", "aaaaaaaa", "aaaa", "aa"]
	OK ["cd"]
	OK ["efcd", "cd"]
	OK ["abcdet2b!end", "cdet2b!", "t2b!", "t2", "b!"]
	OK ["bXcX", "c"]
	OK ["aaa,b,c,d", "c,"]
	OK ["word d e word", "word"]
	OK ["nonnoncapturinggroups"]
	All module tests (14) passed
\end{lstlisting}

\paragraph{Модуль \texttt{matches}}

Тесты функции \verb`matches` проверяют: извлечение всех символов строки, поведение якорей \verb`^` и \verb`$` при глобальном поиске, обработку пустых совпадений (шаблоны \verb`.?` и \verb`.??`), поиск нескольких адресов электронной почты в тексте, и многострочный режим \verb`(?m)` с якорем \verb`$`. Всего 10 тестов.

Результат выполнения модуля \verb`matches` представлен на рисунке~\ref{fig:test_matches}.

\begin{lstlisting}[caption={Результат выполнения модуля matches}, label={fig:test_matches}]
	--- Module "matches" ---
	OK ["a", "b", "c"]
	OK ["a"]
	OK ["bc"]
	OK []
	OK ["a", "b", "c", ""]
	OK ["", "", "", ""]
	OK ["abc@d.ef", "fb@swsu.ru"]
	OK ["<three>", "<cool>", "<words>"]
	OK [""]
	OK ["", "", "", ""]
	All module tests (10) passed
\end{lstlisting}

\paragraph{Модуль \texttt{replace}}

Модуль проверяет корректность функции \verb`re_nreplace` -- замены совпадений в буфере фиксированного размера. Функция \verb`assert_re_nreplace` сравнивает содержимое буфера с эталонной строкой и число выполненных замен с эталонным счётчиком. При \verb`ref_cnt = -1` число замен не проверяется (используется для особых тестов).

Тесты проверяют: удаление символов (пустая строка замены), замену с ограничением числа итераций, использование обратных ссылок \verb`$0`--\verb`$9` и \verb`\n` в строке замены, многострочный режим с якорями \verb`^` и \verb`$`, экранирование метасимволов через замену, и практический сценарий удаления висячих пробелов из кода без нарушения отступов. Два теста помечены как особые (\verb`SPECIAL_TEST`) -- замена на символ при нулевом захвате. Всего 15 тестов.

Результат выполнения модуля \verb`replace` представлен на рисунке~\ref{fig:test_replace}.

\begin{lstlisting}[caption={Результат выполнения модуля replace}, label={fig:test_replace}]
	--- Module "replace" ---
	OK "abc", repl_cnt = 5
	OK "ab c&", repl_cnt = 3
	OK "abc", repl_cnt = 3
	OK "Wikipedia encyclopedia. ", repl_cnt = 4
	OK "abc_def__bla_bla", repl_cnt = 4
	OK "abc_def_bla_bla", repl_cnt = 3
	OK "def", repl_cnt = 1
	OK "{
		int a = 0;
		
		char b = 'a';
		
	}", repl_cnt = 3
	OK "nums: 10, 200, 1040...", repl_cnt = 3
	OK "_-t2b!      b!-", repl_cnt = 1
	OK "li\^te\(r\)al \[match\]\!", repl_cnt = 6
	OK "//a = 1;
	//b = 2;
	//", repl_cnt = 3
	OK "a = 1;", repl_cnt = 1
	OK "a = 1;
	b = 2;
	c = 3;", repl_cnt = 3
	OK "cd
	aa
	//", repl_cnt = 1
	* OK ".t.e.s.t.", repl_cnt = 5
	* OK ".b..c.", repl_cnt = 4
	All module tests (15) passed
\end{lstlisting}

\paragraph{Итоги тестирования}

Итоговая статистика тест-сьюта приведена на рисунке~\ref{fig:test_summary}. Все 237 учитываемых тестов успешно пройдены; 5 особых тестов также дают ожидаемый результат. Суммарное время выполнения составляет 0,022~с.

\begin{lstlisting}[caption={Результат тестирования библиотеки}, label={fig:test_summary}]

-----------------------
All tests (237) passed, elapsed time: 0.022 s
\end{lstlisting}

\subsection{Поисковое приложение}

\subsubsection{Описание классов приложения}

\subsubsubsection{Класс \texttt{TextEditor}}

Класс \verb`TextEditor` является главным окном приложения и наследует \verb`QMainWindow`. Методы класса приведены в таблице~\ref{tab:rp_texteditor_methods}.

\begin{xltabular}{\textwidth}{|C{6cm}|T|}
	\caption{Методы класса \texttt{TextEditor}}
	\label{tab:rp_texteditor_methods} \\
	\hline
	\thead{Метод} & \thead{Описание} \\
	\hline
	\finishhead
	\verb`openFile(path)` & Открывает файл по заданному пути: читает содержимое, определяет тип переноса строк, конвертирует CP1251 в UTF-8 при необходимости, отображает в редакторе \\
	\hline
	\verb`saveToFile(path)` & Сохраняет текст редактора в файл с восстановлением оригинального типа переноса строк (\verb`LF`/\verb`CR`/\verb`CRLF`) \\
	\hline
	\verb`setTabWidth(spaces)` & Устанавливает визуальную ширину табуляции в заданное число пробелов \\
	\hline
	\verb`setZoom(percent)` & Устанавливает масштаб отображения текста; размер шрифта привязывается к ближайшему стандартному значению из списка \verb`ZOOM_SIZES` \\
	\hline
	\verb`detectLineEnding(raw)` & Определяет тип переноса строк в текстовом файле: \verb`CRLF`, \verb`CR` или \verb`LF` \\
	\hline
	\verb`on_action_open_triggered()` & Открывает диалог выбора файла; отклоняет двоичные файлы с сообщением об ошибке \\
	\hline
	\verb`on_action_save_triggered()` & Сохраняет файл по текущему пути; при отсутствии пути (новый файл) вызывает «Сохранить как» \\
	\hline
	\verb`on_action_save_as_triggered()` & Открывает диалог сохранения файла \\
	\hline
	\verb`on_action_find_triggered()` & Показывает диалоговое окно поиска; позиционирует его справа от главного окна или по центру экрана при нехватке места \\
	\hline
	\verb`on_action_search_dir_triggered()` & Открывает диалог выбора директории и запускает \verb`SearchDir` \\
	\hline
	\verb`on_action_zoom_` \verb`in` / \verb`out` / \verb`100` \verb`_triggered()` & Управление масштабом отображения текста в редакторе \\
	\hline
	\verb`onCursorPositionChanged()` & Обновляет индикатор строки и столбца в строке состояния \\
	\hline
	\verb`eventFilter(obj, event)` & Перехватывает \verb`Ctrl+Колёсико` для масштабирования \\
\end{xltabular}

\subsubsubsection{Класс \texttt{SearchDialog}}

Класс \verb`SearchDialog` наследует \verb`QDialog` и реализует немодальный диалог поиска и замены. Методы класса приведены в таблице~\ref{tab:rp_searchdialog_methods}.

\begin{xltabular}{\textwidth}{|C{5cm}|T|}
	\caption{Методы класса \texttt{SearchDialog}}
	\label{tab:rp_searchdialog_methods} \\
	\hline
	\thead{Метод} & \thead{Описание} \\
	\hline
	\finishhead
	\verb`buildPattern()` & Формирует строку шаблона регулярного выражения: добавляет в неё префикс модификаторов \verb`(?Ims)` на основе состояния чекбоксов учёта регистра, многострочного режима и <<\verb`.`>> включает новую строку \\
	\hline
	\verb`buildFindFlags()` & Формирует флаги \verb`QTextDocument::FindFlags` для обычного режима поиска \\
	\hline
	\verb`getSearchScope(text, offset)` & Определяет область поиска: весь текст или только выделенный участок; заполняет \verb`search_text` и \verb`scope_utf8_offset` \\
	\hline
	\verb`validatePattern()` & Компилирует шаблон через \verb`compile_re2`; при ошибке отображает форматированное сообщение с выделением позиции ошибки \\
	\hline
	\verb`collectMatches` \verb`(pattern, text, offset)` & Итеративно вызывает \verb`matchpos` и заполняет список \verb`found_matches` всеми совпадениями в области поиска \\
	\hline
	\verb`highlightAll()` & Формирует список \verb`QTextEdit::ExtraSelection` из \verb`found_matches` и передаёт в \verb`setExtraSelections` \\
	\hline
	\verb`clearHighlights()` & Сбрасывает список совпадений и снимает подсветку \\
	\hline
	\verb`findNext_regex(wrap)` & Поиск следующего совпадения по регулярному выражению с возможным циклическим переходом; переводит курсор редактора на найденное совпадение \\
	\hline
	\verb`findNext_normal(wrap)` & Поиск следующего совпадения в обычном режиме через \verb`QTextDocument::find` с учётом области поиска \\
	\hline
	\verb`utf8ToQcharPos(utf8, pos)` & Конвертирует байтовое смещение UTF-8 в позицию символа Qt \\
	\hline
	\verb`showStatus(html)` / \verb`clearStatus()` & Показывает/скрывает строку состояния в диалоге \\
	\hline
	\verb`on_btn_find_next_clicked()` & Выполняет поиск следующего совпадения \\
	\hline
	\verb`on_btn_count_clicked()` & Подсчитывает число совпадений и отображает в строке состояния \\
	\hline
	\verb`on_btn_replace_clicked()` & Заменяет текущее выделенное совпадение через \verb`re_dupreplace` с \verb`limit=1` \\
	\hline
	\verb`on_btn_replace_all` \verb`_highlight_clicked()` & В режиме замены -- заменяет все совпадения; в режиме поиска -- подсвечивает в тексте все совпадения \\
	\hline
	\verb`on_btn_toggle_replace` \verb`_clicked()` & Переключает видимость панели замены; меняет текст кнопки-переключателя \\
\end{xltabular}

\subsubsubsection{Класс \texttt{SearchDir}}

Класс \verb`SearchDir` наследует \verb`QDialog` и реализует диалоговое окно рекурсивного многопоточного поиска по файлам выбранной директории. Методы класса приведены в таблице~\ref{tab:rp_searchdir_methods}.

\begin{xltabular}{\textwidth}{|C{5cm}|T|}
	\caption{Методы класса \texttt{SearchDir}}
	\label{tab:rp_searchdir_methods} \\
	\hline
	\thead{Метод} & \thead{Описание} \\
	\hline
	\finishhead
	\verb`startSearch()` & Компилирует шаблон, формирует список фильтров расширений, создаёт \verb`QThread` и \verb`SearchController`, соединяет сигналы и запускает поток \\
	\hline
	\verb`stopSearch()` & Устанавливает флаг остановки в \verb`SearchController` через \verb`requestStop()` \\
	\hline
	\verb`cleanupThread()` & Завершает и удаляет рабочий поток после окончания поиска \\
	\hline
	\verb`onMatchFound(result)` & Добавляет дочерний узел совпадения в дерево результатов; при необходимости создаёт родительский узел файла \\
	\hline
	\verb`onFileProcessed` \verb`(processed, total)` & Обновляет прогресс-бар \\
	\hline
	\verb`onSearchFinished` \verb`(files, matches, stopped)` & Скрывает прогресс-бар, обновляет счётчики в узлах файлов, отображает итоговый статус \\
	\hline
	\verb`onTreeItemDoubleClicked` \verb`(item, col)` & Открывает файл в редакторе через \verb`TextEditor::openFile` и переводит курсор на строку совпадения \\
\end{xltabular}

\subsubsubsection{Классы \texttt{SearchController} и \texttt{FileSearchWorker}}

Класс \verb`SearchController` наследует \verb`QObject` и выполняется в отдельном потоке. Он обходит директорию через \verb`QDirIterator` и создаёт объекты \verb`FileSearchWorker` -- по одному на каждый файл. Воркеры помещаются в локальный \verb`QThreadPool`; при многопоточном режиме (\verb`chk_multithread`) пул использует все доступные ядра процессора, при однопоточном -- одно. Методы классов приведены в таблице~\ref{tab:rp_worker_methods}.

\begin{xltabular}{\textwidth}{|C{5cm}|T|}
	\caption{Методы классов \texttt{SearchController} и \texttt{FileSearchWorker}}
	\label{tab:rp_worker_methods} \\
	\hline
	\thead{Метод} & \thead{Описание} \\
	\hline
	\finishhead
	\verb`SearchController::setup(...)` & Инициализирует параметры поиска: директорию, фильтры по расширениям, флаги режима поиска, скомпилированный объект регулярного выражения \\
	\hline
	\verb`SearchController::run()` & Вычисляет суммарный размер файлов, обходит директорию, создаёт и запускает воркеры, ожидает их завершения через \verb`QEventLoop` \\
	\hline
	\verb`SearchController` \verb`::requestStop()` & Устанавливает атомарный флаг остановки \verb`stop_flag` \\
	\hline
	\verb`FileSearchWorker::run()` & Читает файл, проверяет на двоичность по наличию нулевых байтов в его первых 8~КБ, конвертирует CP1251 в UTF-8, итеративно ищет совпадения и испускает сигнал \verb`matchFound` для каждого \\
	\hline
	\verb`FileSearchWorker` \verb`::isBinary(data)` & Проверяет буфер на наличие нулевых байтов \\
	\hline
	\verb`FileSearchWorker` \verb`::extractContext(...)` & Извлекает фрагмент текста вокруг совпадения для отображения в дереве результатов \\
\end{xltabular}

\subsubsection{Тестирование поискового приложения}

Тестирование поискового приложения проводится вручную по сценариям, охватывающим позитивные, негативные и граничные случаи. Для каждого сценария приведены входные данные, ожидаемый результат и скриншот с результатом.

\subsubsubsection{Сценарий 1. Открытие текстового файла}

Входные данные: текстовый файл в кодировке UTF-8 с содержимым:

\begin{lstlisting}[language={}]
	Строка первая
	Строка вторая
	Строка третья
\end{lstlisting}

Ожидаемый результат: файл открывается, содержимое отображается в редакторе, в строке состояния показаны «Стр: 1  Стб: 1», тип переноса строк «LF» и масштаб «100\%», заголовок окна содержит имя файла.

Результат тестирования представлен на рисунке~\ref{fig:test_app_open}.

\begin{figure}[H]
	\centering
	\includegraphics[width=0.95\linewidth]{test1.png}
	\caption{Открытие текстового файла}
	\label{fig:test_app_open}      
\end{figure}

\subsubsubsection{Сценарий 2. Открытие файла в кодировке CP1251}

Входные данные: файл \verb`.txt` в кодировке Windows-1251 с кириллическим текстом.

Ожидаемый результат: кириллица отображается корректно (автоматическая конвертация в UTF-8); в строке состояния отображается тип переноса строк (CRLF).

Результат тестирования представлен на рисунке~\ref{fig:test_app_cp1251}.

\begin{figure}[H]
	\centering
	\includegraphics[width=0.95\linewidth]{test2.png}
	\caption{Открытие файла в кодировке CP1251}
	\label{fig:test_app_cp1251}      
\end{figure}

\subsubsubsection{Сценарий 3. Попытка открытия двоичного файла}

Входные данные: двоичный файл (например, \verb`.exe`, \verb`.png` или \verb`.bin`).

Ожидаемый результат: появляется диалог с сообщением «Файл является двоичным и не может быть открыт в текстовом редакторе»; содержимое редактора не изменяется.

Результат тестирования представлен на рисунке~\ref{fig:test_app_binary}.

\begin{figure}[H]
	\centering
	\includegraphics[width=0.95\linewidth]{test3.png}
	\caption{Попытка открытия двоичного файла}
	\label{fig:test_app_binary}      
\end{figure}

\subsubsubsection{Сценарий 4. Поиск по регулярному выражению с подсветкой}

Входные данные: открытый файл со следующим содержимым:

\begin{lstlisting}[language={}]
	email: user@example.com
	phone: +79161234567
	email: admin@test.org
\end{lstlisting}

Шаблон: \verb`[\w.+-]+@[\w-]+\.[\w.]+`, режим -- регулярное выражение, нажата кнопка «Подсветить».

Ожидаемый результат: оба адреса электронной почты подсвечены; в строке состояния отображается «Всего 2 совпадений в файле».

Результат тестирования представлен на рисунке~\ref{fig:test_app_highlight}.

\begin{figure}[H]
	\centering
	\includegraphics[width=0.95\linewidth]{test4.png}
	\caption{Поиск по регулярному выражению с подсветкой}
	\label{fig:test_app_highlight}      
\end{figure}

\subsubsubsection{Сценарий 5. Поиск с некорректным шаблоном}

Входные данные: шаблон \verb`.[\w-z]`, режим -- регулярное выражение.

Ожидаемый результат: в строке состояния диалога отображается сообщение об ошибке красным цветом с подчёркиванием символа, вызвавшего ошибку (дефис, обозначающий диапазон символов в наборе символов); поиск не выполняется.

Результат тестирования представлен на рисунке~\ref{fig:test_app_error}.

\begin{figure}[H]
	\centering
	\includegraphics[width=0.7\linewidth]{test5.png}
	\caption{Попытка поиска с некорректным шаблоном}
	\label{fig:test_app_error}      
\end{figure}

\subsubsubsection{Сценарий 6. Некорректный шаблон с ошибкой в недостающем символе}

Входные данные: шаблон \verb`(abc` с ошибкой в недостающем символе (после текста шаблона), режим -- регулярное выражение.

Ожидаемый результат: в строке состояния диалога отображается сообщение об ошибке красным цветом с подчёркиванием места возникновения ошибки (нуль-терминатор шаблона -- отсутствует закрывающая скобка); поиск не выполняется.

Результат тестирования представлен на рисунке~\ref{fig:test_app_error_edge}.

\begin{figure}[H]
	\centering
	\includegraphics[width=0.7\linewidth]{test6.png}
	\caption{Шаблон с ошибкой в недостающем символе}
	\label{fig:test_app_error_edge}      
\end{figure}

\subsubsubsection{Сценарий 7. Замена всех совпадений с обратными ссылками}

Входные данные: открыт файл:
\\

\begin{lstlisting}[language={}]
	Дата: 15.03.2023
	Срок: 31.12.2025
\end{lstlisting}

Шаблон: \verb`(\d{2})\.(\d{2})\.(\d{4})`, строка замены: \verb`$3-$2-$1`, нажата кнопка «Заменить всё».

Ожидаемый результат: содержимое редактора принимает вид:

\begin{lstlisting}[language={}]
	Дата: 2023-03-15
	Срок: 2025-12-31
\end{lstlisting}

В строке состояния отображается «Выполнено 2 замен».

Результат тестирования представлен на рисунке~\ref{fig:test_app_replace}.

\begin{figure}[H]
	\centering
	\includegraphics[width=0.95\linewidth]{test7.png}
	\caption{Содержимое редактора после выполнения замены}
	\label{fig:test_app_replace}      
\end{figure}

\subsubsubsection{Сценарий 8. Поиск с циклическим переходом}

Входные данные: файл с несколькими совпадениями, курсор на последнем совпадении, включена галочка «Зациклить поиск».

Ожидаемый результат: при нажатии «Найти далее» поиск переходит к первому совпадению; в строке состояния отображается «Достигнут конец, поиск продолжен с начала файла».

Результат тестирования представлен на рисунке~\ref{fig:test_app_wrap}.

\begin{figure}[H]
	\centering
	\includegraphics[width=0.95\linewidth]{test8.png}
	\caption{Выделение после циклического перехода}
	\label{fig:test_app_wrap}      
\end{figure}

\subsubsubsection{Сценарий 9. Поиск в выделенной области}

Входные данные: выделен фрагмент текста, включена галочка «Только в выделении», задан шаблон.

Ожидаемый результат: подсветка и подсчёт совпадений выполняются только внутри выделенного фрагмента; символы за пределами выделения не затрагиваются.

Результат тестирования представлен на рисунке~\ref{fig:test_app_scope}.

\begin{figure}[H]
	\centering
	\includegraphics[width=0.95\linewidth]{test9.png}
	\caption{Поиск в выделении}
	\label{fig:test_app_scope}      
\end{figure}

\subsubsubsection{Сценарий 10. Рекурсивный поиск по директории}

Входные данные: директория с файлами:

\begin{lstlisting}[language={}]
	file1.txt  (содержит "TODO: fix bug")
	subdir/file2.cpp  (содержит "TODO: refactor")
	subdir/image.png  (двоичный файл)
\end{lstlisting}

Шаблон: \verb`TODO:.*`, фильтр по расширениям: \verb`*`, рекурсивный режим включён.

Ожидаемый результат: в дереве результатов отображаются два файла с одним совпадением каждый; файл \verb`image.png` пропущен как двоичный; при двойном щелчке на совпадении файл открывается в редакторе с переходом на нужную строку.

Результат тестирования представлен на рисунке~\ref{fig:test_app_searchdir}.

\begin{figure}[H]
	\centering
	\includegraphics[width=0.95\linewidth]{test10.png}
	\caption{Результаты рекурсивного поиска}
	\label{fig:test_app_searchdir}      
\end{figure}

\subsubsubsection{Сценарий 11. Прерывание поиска по директории}

Входные данные: директория с большим числом файлов (более 1000), поиск запущен.

Ожидаемый результат: при нажатии кнопки «Прервать поиск» поиск останавливается; кнопка возвращается в состояние «Найти»; в строке состояния отображается «Поиск прерван»; уже найденные результаты остаются в дереве.

Результат тестирования представлен на рисунке~\ref{fig:test_app_stop}.

\begin{figure}[H]
	\centering
	\includegraphics[width=0.95\linewidth]{test11.png}
	\caption{Прерывание поиска по директории}
	\label{fig:test_app_stop}      
\end{figure}

\subsubsubsection{Сценарий 12. Граничный случай: совпадение нулевой длины}

Входные данные: шаблон соответствует захвату нулевой длины (например, \verb`.??`).

Ожидаемый результат: в строке состояния отображается «Совпадение нулевой длины -- выделен ближайший символ»; ошибок не возникает.

Результат тестирования представлен на рисунке~\ref{fig:empty_capture}.

\begin{figure}[H]
	\centering
	\includegraphics[width=0.95\linewidth]{test12.png}
	\caption{Обработка совпадения нулевой длины}
	\label{fig:empty_capture}      
\end{figure}

\subsubsubsection{Сценарий 13. Граничный случай: шаблон совпадает со всем текстом}

Входные данные: шаблон \verb`(?s).+` применяется к многострочному тексту, "<Многострочный режим"> отключён, режим "<\textquotedbl .\textquotedbl \ -- новая строка"> включается через модификатор в составе строки шаблона.

Ожидаемый результат: всё содержимое редактора подсвечивается как одно совпадение; подсчёт возвращает 1.

Результат тестирования представлен на рисунке~\ref{fig:test_app_dotall}.

\begin{figure}[H]
	\centering
	\includegraphics[width=0.95\linewidth]{test13.png}
	\caption{Обработка шаблона, захватывающего весь многострочный текст}
	\label{fig:test_app_dotall}      
\end{figure}